
\subsection*{2.5}
\phantomsection
\label{sga3:VIIB:2.5}
\addcontentsline{toc}{subsection}{2.5. Étale formal groups}

A formal \(k\)-group is said to be \emph{étale} if its underlying
formal variety is étale \textup{(cf.\ \SGARefLink{sga3:VIIB:1.6}{1.6})}. Such formal groups have
a very simple structure. Suppose that \(k\) is local, let \(\kappa\)
be its residue field, let \(\kappa_s\) be a separable closure of
\(\kappa\), and let
\[
  \Gamma=\operatorname{Gal}(\kappa_s/\kappa)
\]
be the topological Galois group. A \emph{\(\Gamma\)-set} means a set
\(E\) with a continuous action of \(\Gamma\); equivalently, the
stabilizer of every element \(x\in E\) is an open subgroup of
\(\Gamma\).

For every formal \(k\)-variety \(X\), put
\[
  X(\kappa_s)
  =
  \varinjlim_{\ell}X(\ell),
\]
where \(\ell\) ranges over the finite extensions of \(\kappa\)
contained in \(\kappa_s\).\footnote{%
\label{sga3:VIIB:nde:95}If
\([\kappa_s:\kappa]=\infty\), then \(\kappa_s\) is not a profinite
\(\kappa\)-algebra, so the preceding notation is an abuse. The
editors also supplied the details that follow.}
The group \(\Gamma\) acts continuously on every \(X(\ell)\), and
hence on \(X(\kappa_s)\). Moreover, if
\[
  X_\kappa=X\widehat\otimes_k\kappa
  \qquad\text{\textup{(cf.\ 1.6.C)}},
\]
then \(X(\ell)=X_\kappa(\ell)\) for every \(\ell\), and therefore
\[
  X(\kappa_s)=X_\kappa(\kappa_s).
\]

Now suppose that \(X\) is étale over \(k\). Then \(X_\kappa\) is the
formal \(\kappa\)-variety that is the disjoint union of the
\(\operatorname{Spec}\kappa(x)\), for \(x\in X\). If
\(X'_\kappa\) denotes the \(\kappa\)-scheme given by the same
disjoint union, then
\[
\begin{aligned}
  X_\kappa(\kappa_s)
  &=
  X'_\kappa(\kappa_s)
  \\
  &=
  \operatorname{Hom}_{\operatorname{Sch}/\kappa}
  \bigl(\operatorname{Spec}\kappa_s,X'_\kappa\bigr).
\end{aligned}
\]

Let \(\mathcal C\) be the full subcategory
\[
  \operatorname{Sch}^{\mathrm{\acute{e}t}}/\kappa
  \subset
  \operatorname{Sch}/\kappa
\]
of étale \(\kappa\)-schemes. It is known that
\[
  X'\longmapsto X'(\kappa_s)
\]
is an equivalence from \(\mathcal C\) onto the category
\(\mathcal C'\) of \(\Gamma\)-sets
\textup{(cf.\ SGA 1, V, §§7--8, or \([DG70]\), §I.4, 6.4)}.
It therefore induces an equivalence between the category of group
objects in \(\mathcal C\) and the category of group objects in
\(\mathcal C'\). A group object in \(\mathcal C'\) is simply an
abstract group endowed with a continuous action of \(\Gamma\) by
group automorphisms; we call this a \emph{\(\Gamma\)-group}.

Taking into account the equivalences of \SGARefLink{sga3:VIIB:1.6.E}{1.6.E},
\[
  \operatorname{Vaf}^{\mathrm{\acute{e}t}}/k
  \xrightarrow{\ \sim\ }
  \operatorname{Vaf}^{\mathrm{\acute{e}t}}/\kappa
  \xrightarrow{\ \sim\ }
  \operatorname{Sch}^{\mathrm{\acute{e}t}}/\kappa,
\]
we obtain the following.

\medskip
\noindent\textbf{Proposition.}
\phantomsection
\label{sga3:VIIB:2.5:proposition}
{\itshape Let \(k\) be a local pseudocompact ring, \(\kappa\) its
residue field, \(\kappa_s\) a separable closure of \(\kappa\), and
\(\Gamma=\operatorname{Gal}(\kappa_s/\kappa)\).
\begin{enumerate}
\item[\textup{(i)}]
The functor \(X\mapsto X(\kappa_s)\) is an equivalence from the
category of étale formal \(k\)-varieties onto the category of
\(\Gamma\)-sets.

\item[\textup{(ii)}]
It induces an equivalence from the category of étale formal
\(k\)-groups onto the category of \(\Gamma\)-groups.
\end{enumerate}}

\medskip
\noindent\textbf{Remark 2.5.A.}
\phantomsection
\label{sga3:VIIB:2.5.A}
\footnote{\label{sga3:VIIB:nde:96}The editors added this remark.}
Let \(k\) be a field, let \(k_s\) be a separable closure of \(k\),
let \(G\) be an étale formal \(k\)-group, and put
\[
  M=G(k_s).
\]
Choose a set \(X\) of representatives for the orbits of
\(\Gamma=\operatorname{Gal}(k_s/k)\) in \(M\). For each \(x\in X\),
let \(\Gamma_x\) be its stabilizer, an open subgroup of finite index
in \(\Gamma\), and put
\[
  L_x=k_s^{\Gamma_x}.
\]
Then \(L_x/k\) has degree \([\Gamma:\Gamma_x]\)
\textup{(see, for example, \([BAlg]\), V, §10.10)}.

By the preceding equivalence of categories, the affine algebra
\(\mathcal A(G)\) of \(G\) is the product of the \(L_x\), endowed
with the product topology. The \(L_x\) are therefore exactly the
simple quotients
\[
  \mathcal A(G)/\mathfrak m,
\]
where \(\mathfrak m\) ranges over the open maximal ideals of
\(\mathcal A(G)\). Since these quotients correspond by duality to
the simple subcoalgebras of \(\mathbf H(G)\), the coalgebra
\(\mathbf H(G)\) is pointed---that is,
\[
  \dim_kC=1
\]
for every simple subcoalgebra \(C\)---if and only if \(L_x=k\) for
every \(x\). In that case, \(\mathcal A(G)\) is the topological
algebra \(k^M\), so \(\mathbf H(G)\) is the group algebra \(kM\), and
\[
  M
  =
  \bigl\{
    x\in\mathbf H(G)
    \mathrel{\big|}
    \varepsilon(x)=1,\ \Delta(x)=x\otimes x
  \bigr\}
  =
  G(k).
\]
